\subsection{标准群中的幂}

保留第 4 号的记号。

命题 8. 设 \(n \in \mathbf{Z}\)\(h_n\)，h＿n 是从 G 到 G 的映射 \(x \mapsto x^n\)\(G\)\(G\)\(a\)\(A\)\(m\)。设 a 是包含于 m 中的 A 的非零理想，且 \(n \notin a\)\(h_n|G(a)\)\(G(a)\)\(G(na)\)。则 h＿n 在 G(a) 上的限制是从解析流形 G(a) 到解析流形 G(na) 的同构。

由标准群的定义，\(h_n\) 在整个 G 上等\(G\)于一个系数在 A＾r 中的整级数。由\(A^r\) § 5，公式 (4)，此级数形如

\[
h _ {n} (x) = n x + \sum_ {| \alpha | \geqslant 2} a _ {\alpha} x ^ {\alpha}.
\]

因此，对于 \(x \in G\)，

\[
\begin{array}{r l} h _ {n} (n x) & = n ^ {2} \Big (x + \sum_ {| \alpha | \geqslant 2} a _ {\alpha} n ^ {| \alpha | - 2} x ^ {\alpha} \Big) \\ & = n ^ {2} \mathrm{S} (x) \end{array}
\]

其中我们记 \(\mathrm{S}(x)=x+\sum_{|\alpha| \geq 2} a_{\alpha} n^{|\alpha|-2} x^{\alpha}\)。此级数 \(\mathrm{S}(x)\) 定义了一个从 G 到 G 的解析映射，仍记为 S。由代数，第 IV 章，§6，命题 8，存在一个系数在  中的整级数 \(S'\)，使得 \(A^r\)\(\mathrm{S}'(\mathrm{S}(\mathrm{X}))=\mathrm{S}(\mathrm{S}'(\mathrm{X}))=\mathrm{X}\)。因此 S 是从解析流形 G 到自身的同构，并且对于 A 的任意包含于 m 中的非零理想 b，都有 \(\mathrm{S}(\mathrm{G}(\mathrm{b}))\subset\mathrm{G}(\mathrm{b}),\mathrm{S}'(\mathrm{G}(\mathrm{b}))\subset\mathrm{G}(\mathrm{b})\)、，从而

\[
\mathrm{S} (\mathrm{G} (\mathfrak {b})) = \mathrm{G} (\mathfrak {b}).
\]

由于 \(h_n(y) = n^2\mathrm{S}\left(\frac{1}{n} y\right)\) 对 \(y \in n\mathrm{G}\) 成立，可见 h＿n 在 \(h_n|n\mathrm{G}(\mathfrak{b})\)\(n\mathrm{G}(\mathfrak{b})\) 上的限制是从解析流形  到解析流形 \(n^2\mathrm{G}(\mathfrak{b})\) 的同构。但是，由于 \(n \notin \mathfrak{a}\)，对所有 \(|n| > |\lambda|\) 都有 \(\lambda \in \mathfrak{a}\)，从而 \(n^{-1}\mathfrak{a} \subset \mathfrak{m}\)，因此 \(\mathfrak{a}\) 具有形式 \(n\mathfrak{b}\)，其中 \(\mathfrak{b}\) 是包含于 \(\mathrm{A}\)\(\mathfrak{m}\) 中的 A 的非零理想。

推论。若 n 在 A 中可逆，则 h＿n 是从解析流形 G 到自身的同构。对于 A 的每个包含于 \(n\)\(h_n\)\(\mathfrak{a}\) 中的非零理想 \(\mathfrak{m}\)，都有 \(h_n(\mathrm{G}(\mathfrak{a})) = \mathrm{G}(\mathfrak{a})\)。对所有 \(x \in \mathrm{G}\)，都有 \(\omega(x^n) = \omega(x)\)。

这立即由命题 8 得到。

命题 9. 假设 \(p \neq 0\)。

\begin{enumerate}
\def\labelenumi{(\roman{enumi})}
\item
  设 \(\mathfrak{a},\mathfrak{b}\) 是 \(\mathbf{A}\) 的非零理想，满足 \(\mathfrak{b} \subset \mathfrak{a} \subset \mathfrak{m}\)。在群 \(\mathrm{G}(\mathfrak{a}) / \mathrm{G}(\mathfrak{b})\)\(p\) 中，每个元素的阶都是 p 的幂。
\item
  假设 \(v(p) = 1\)。若 \(x \in G\) 满足 \(\omega(x) > \frac{1}{p - 1}\)，则
\end{enumerate}

\[
\omega (x ^ {p}) = \omega (x) + 1.
\]

由 § 5，公式 (4)，对所有 \(x \in G\)，

\[
x ^ {p} = p x + \sum_ {| \alpha | \geqslant 2} c _ {\alpha} x ^ {\alpha}
\]

即使为了证明 (i)，也可以假设 \(c_{\alpha} \in A^{r}\)\(\alpha\)\(v(p) = 1\)。于是若 \(\omega(x) \geqslant 1\)，便有 \(\omega(x^{p}) \geqslant \omega(x) + 1\)，从而当 n 趋于 \(\omega(x^{p^{n}})\) 时 \(+\infty\)\(n\) 趋于 \(+\infty\)；这证明了 (i)。由于 \(\binom{p}{i}\)\(p\) 可被 p 整除，对于

\(1 \leqslant i \leqslant p - 1\)\(c_{\alpha} \in pA^{r}\)，§ 5，第 3 号，命题 2 说明，对于

\[
2 \leqslant | \alpha | \leqslant p - 1
\]

从而

\[
\omega \left(c _ {\alpha} x ^ {\alpha}\right) > \omega (p x) = \omega (x) + 1 \quad \text { for } 2 \leqslant | \alpha | \leqslant p - 1.
\]

另一方面，若 \(|\alpha| \geqslant p\)，则 \(\omega(c_{\alpha}x^{\alpha}) \geqslant p\omega(x)\)，而当 \(p\omega(x) > \omega(x) + 1\) 时，\(\omega(x) > \frac{1}{p - 1}\)。这证明了 (ii)。

